\section{The upper bound for arbitrary linear restrictions}\label{sec:upper}

A \emph{slice in mode $i$} is a tensor of the form $a\otimes B$,
with $a$ in the $i$th factor and $B$ in the tensor product of the other
factors, placed in their original order. The \emph{slice rank} $\sr(T)$
is the least number of slices summing to $T$; their distinguished modes
may differ. Applying linear maps independently to the factors takes
a slice to a slice or zero. We first prove the diagonal value needed
to bound ordinary subrank. This fact and the entropy upper-bound method
have precedents in \citet{TaoSawin2016}; the required arguments follow.

\begin{lemma}\label{lem:product-dimension}
For subspaces $U,V\subseteq\C^q$, let $U*V$ be the linear span of
their coordinatewise products. Then
\begin{equation}\label{eq:product-dimension}
 \dim(U*V)\geq\dim U+\dim V-q.
\end{equation}
Consequently, for any $\ell\geq1$ and $U_1,\ldots,U_\ell\subseteq\C^q$,
\begin{equation}\label{eq:iterated-product}
 \dim\operatorname{span}\{u_1*\cdots*u_\ell:u_i\in U_i\}
    \geq q-\sum_{i=1}^{\ell}\codim U_i.
\end{equation}
\end{lemma}
\begin{proof}
Let $A$ be the union of the coordinate supports of vectors in $U$.
There is $u\in U$ nonzero in every coordinate of $A$: each such
coordinate is a nonzero linear function on $U$, and a finite union of
their proper kernels cannot cover a vector space over the infinite
field $\C$. Equivalently, the product of these linear functions is
a nonzero polynomial and is nonzero somewhere. If $U=0$, take $A$
empty and $u=0$.

Coordinatewise multiplication by $u$ on $\C^q$ has kernel of dimension
$q-|A|$. Since $\dim U\leq|A|$, its restriction to $V$ gives
\[
 \dim(U*V)\geq\dim(u*V)
     \geq\dim V-(q-|A|)\geq\dim V+\dim U-q.
\]
Iterate this inequality, using bilinearity of coordinatewise
multiplication to identify the successive spans. The case $\ell=1$
of \eqref{eq:iterated-product} is an equality.
\end{proof}

\begin{lemma}\label{lem:diagonal-slice}
For every $k\geq2$ and $q\geq1$, $\sr(\Delta_q)=q$.
In particular, $Q(T)\leq\sr(T)$ for every $k$-tensor $T$ over $\C$.
\end{lemma}
\begin{proof}
Consider any slice decomposition of $\Delta_q$. For each $i$, let
$r_i$ be the number of its slices in mode $i$, and let $U_i\subseteq\C^q$
annihilate their distinguished coefficient vectors under the bilinear
pairing $\sum_j x_jy_j$. Then $\codim U_i\leq r_i$. Evaluating the
decomposition on $U_1\times\cdots\times U_k$ kills every slice, so
\[
 \sum_{j=1}^q u_{1,j}\cdots u_{k,j}=0
       \qquad (u_i\in U_i).
\]
Thus the span of $U_1*\cdots*U_{k-1}$ lies in $U_k^\perp$, whose
dimension is $\codim U_k$. By \eqref{eq:iterated-product},
\[
 q-\sum_{i<k}\codim U_i\leq\codim U_k,
       \qquad q\leq\sum_i\codim U_i\leq\sum_i r_i.
\]
The defining expression for $\Delta_q$ is a sum of $q$ slices, proving
equality. Monotonicity of slice rank under independent linear maps
now gives $q=\sr(\Delta_q)\leq\sr(T)$ for every diagonal restriction
of $T$, and hence $Q(T)\leq\sr(T)$.
\end{proof}

\begin{proposition}\label{prop:upper}
For all $N\geq1$,
\[
 Q(D_{\mathbf d}^{\boxtimes N})\leq
 \sr(D_{\mathbf d}^{\boxtimes N})\leq k\exp(N\HH(p)).
\]
\end{proposition}
\begin{proof}
For a local word $w=(w_1,\ldots,w_N)\in\{1,\ldots,n\}^N$, set
\[
 s(w)=-\sum_{t=1}^N\log p_{w_t},\qquad
 P(w)=\prod_{t=1}^N p_{w_t}=e^{-s(w)},\qquad
 L_N=\{w:s(w)\leq N\HH(p)\}.
\]
All $p_a$ are positive. Since $\sum_wP(w)=1$ and every $w\in L_N$
has $P(w)\geq e^{-N\HH(p)}$, we have
\begin{equation}\label{eq:score-count}
 |L_N|\leq e^{N\HH(p)}.
\end{equation}
For every support edge $(w^{(1)},\ldots,w^{(k)})$ of the full tensor
power, each column has multiplicities $\mathbf d$. Therefore
\[
 \sum_{i=1}^k s(w^{(i)})
    =-N\sum_{a=1}^n d_a\log p_a=Nk\HH(p).
\]
At least one local word belongs to $L_N$. Assign the edge to the least
such mode. For each $i$ and $w\in L_N$, collect all assigned terms
whose $i$th word is $w$ into one slice with distinguished factor $e_w$.
This partitions the entire support and yields at most $k|L_N|$ slices.
Apply Lemma~\ref{lem:diagonal-slice} and \eqref{eq:score-count}.
\end{proof}

Taking $N$th roots in Proposition~\ref{prop:upper} and using
Lemma~\ref{lem:limits} gives $\Qt(D_{\mathbf d})\leq e^{\HH(p)}$.
Together with the coordinate-restriction lower bound
\eqref{eq:monomial-lower} and $\Qmt\leq\Qt$, this proves
Theorem~\ref{thm:main} for $k\geq3$. The two-factor case was handled
in Section~\ref{sec:limits}. Every restriction in the argument retains
the same $k$ factors.
